% Transcription — fonds Alexandre Grothendieck, shelfmark 12, batch 7 (pages 121-140).
% Established from the facsimile archives/batches/12/batch-07.pdf (a mirror of
% https://grothendieck.umontpellier.fr/12.pdf), per the skill
% transcribe-grothendieck. The batch file carries no cover sheet: PDF page N
% is page 120+N, as the Montpellier footer printed on each PDF page confirms,
% and the archivists' pencil numbers bottom-left agree with that position
% wherever they are legible (130, 132, 133, 134, 135, 136, 137, 138, 139, 140).
%
% Pass: Opus 5.5 (claude-opus-5-5), 2026-09-23 — first pass, unchecked against the pages by a human.
%
% What the batch is.
%   121      Handwritten, item « 2) » of a list whose item 1 is not on this
%            leaf: F a
%            sheaf R^d φ_*(Q_ℓ) coming from a square X' → Y' over X → Y;
%            the Leray spectral sequence of fφ degenerates « par Deligne »,
%            so R^i f_*(F) is a piece of R^{i+j}(fφ)_*(Q_ℓ); then the other
%            spectral sequence through Y' and bounds on weight and level.
%   122      One leaf of a French typescript, corrected by hand: item « 4) »
%            of a report on Deligne's work — the cohomological triviality
%            (degeneration of the Leray spectral sequence) of a smooth
%            projective morphism, proved « modulo le théorème de Lefschetz »;
%            Blanchard's thesis as the special case of trivial local
%            system; the consequences for Riemann-Weil with twisted
%            coefficients and for the « théorie des motifs ». He replaces
%            « rédacteur » by « rapporteur » once, not the other time.
%            Breaks off mid-sentence (« qui se »).
%   124-136  (even pages) One handwritten run under the underlined title
%            « Divers types de Q_ℓ-faisceaux » (p. 124, first word doubtful):
%            the definition of an admissible Q_ℓ-sheaf of type ≤ d (locally
%            a finite filtration whose graded pieces are, or are summands of,
%            R^n g_*(Q_ℓ) for g projective smooth with connected fibres of
%            dim ≤ d); Lemme 1 (thick abelian subcategory), Lemme 2 (inverse
%            images, « trivial »), Lemme 3 (p. 126: R^i f_! of a type-≤d sheaf
%            under f of relative dim ≤ d' is of type ≤ d+d'; the « 2 » is
%            struck and made « 3 »), its proof by induction on d+d', dévissage
%            to R^p f_!(R^q g_*(Q_ℓ)), a factorisation d'Artin reducing to
%            dim U = 1, short exact sequences, and resolution of
%            singularities V = V̂ − Z (pp. 128-130). Then « Proposition 1 »
%            (p. 132): for i : U → X open, X excellent of dim ≤ d, F
%            d'-admissible, R^n i_*(F) is (d+d'-1)-admissible (the index
%            doubtful); the proof by induction and noetherian induction,
%            reduction to U regular, resolution to a normal-crossings
%            complement, purity when F = Q_ℓ, compactification of V (p. 134),
%            and on p. 136 the spectral sequence with E_2^{pq} = 0 for
%            q ∉ [0,2] (a sketch of the E_2 term), E_2^{p2} ≃ Q_ℓ(-1) ⊗
%            E_2^{p0}, and the case of E_2^{p1} left « à traiter…!! ».
%            These are the rectos of leaves whose other sides are the odd
%            pages 125-137 — a Bourbaki typescript (n° 383, groups generated
%            by reflections), unmarked by him — skipped.
%   123      Blank. Absent.
%   138      A cover in his hand, boxed: « Motifs lisses versus Isomotifs
%            lisses ». No \page; its words head the next run.
%   139-140  Fair copy on different paper: over S normal integral with
%            generic point η, k(η) = K, a smooth isomotif (T_ℓ(M) unramified
%            over S[1/ℓ]); a smooth motif as a triple (M, {E_ℓ}, {α_ℓ}) with
%            E_ℓ ℓ-adic sheaves on S_ℓ and α_ℓ : E_ℓ ⊗ Q_ℓ ≃ T_ℓ(M), almost
%            all integral; morphisms (u, {u_ℓ}) and the commuting square; the
%            N.B. that motifs over S → motifs over K is fully faithful, the
%            category abelian, motifs → isomotifs exact and essentially
%            surjective with kernel the finite commutative groups étale
%            over each S[1/ℓ]; ⊗-structure; R^i f_*(Z_X) for X projective
%            smooth over S; and « Relations avec les schémas en groupes
%            commutatifs propres et lisses »: the contravariant functor
%            A ↦ (R^1 f_{A°*}(Q), Hom(_{ℓ^∞}A, Q_ℓ/Z_ℓ), u), closing on a
%            N.B. about A/A° finite étale.
%
% Notation fixed for the batch. Blackboard \mathbb{Q}_{\ell}, \mathbb{Z}_{\ell}
% for his double-struck Q, Z; his double-struck R of derived functors (pp.
% 132-136) is \mathbf{R}, the plain R^{n} elsewhere as he writes it; his
% R^{i}_{!}f for what is now R^{i}f_{!} is kept; Q^{*} for the cone on p. 132;
% \mathfrak{P}(A) for the functor on p. 140 (the capital is a reading);
% ${}_{\ell^{\infty}}A$ for ℓ-primary torsion, index on the left as he writes it.
%
% Legibility. Pages 122 (typescript), 139 and 140 read nearly through. Pages
% 124-136 are a fast draft whose formulas carry and whose connective prose
% often does not; 128, 132 and the lower half of 136 are the hardest.
\documentclass[11pt,a4paper]{article}
% Le préambule commun à toutes les transcriptions du fonds Grothendieck.
%
% Il tient en une idée : **l'incertitude se compose, elle ne se résout pas.**
% Une transcription qui lisse un mot douteux a détruit la seule chose pour
% laquelle on la faisait — un lecteur doit pouvoir savoir, à chaque endroit,
% ce qui était sur le feuillet et ce qui a été ajouté. Les six macros qui
% suivent sont donc l'appareil critique, et non de la décoration.
%
% Les mêmes commandes sont reconnues par `scripts/render.mjs`, qui produit la
% vue de lecture du site. Une macro ajoutée ici doit l'être aussi là-bas,
% faute de quoi le rendu échouera bruyamment — ce qui est voulu : un rendu
% silencieusement faux est pire qu'un rendu absent.

\ProvidesPackage{grothendieck}[2026/08/08 Transcriptions du fonds Grothendieck]

\usepackage{amsmath,amssymb,amsthm}
\usepackage{mathrsfs}
\usepackage{stmaryrd}
% Les diagrammes commutatifs sont partout dans ce fonds : c'est la langue
% même de la géométrie algébrique de Grothendieck, et une transcription qui
% les rendrait en prose ne transcrirait rien.
\usepackage{tikz-cd}
% Français par défaut — c'est la langue des feuillets — mais l'anglais est
% chargé aussi : la traduction et le résumé sont des documents anglais, et la
% typographie française (espaces avant les deux-points) y serait une faute.
% Ces documents basculent par \selectlanguage{english} en tête de corps.
\usepackage[english,french]{babel}
\usepackage{fontspec}
\usepackage{geometry}
\geometry{a4paper,margin=25mm,marginparwidth=28mm}
\usepackage{marginnote}
\usepackage{xcolor}
% ulem, not soul. `soul`'s \st and \ul are built on a character-by-character
% scan that XeLaTeX's font machinery breaks: the document fails with an
% undefined control sequence rather than degrading. ulem does the same two jobs
% and is Unicode-engine safe. `normalem` stops it hijacking \emph, which the
% transcriptions use for Grothendieck's own emphasis.
\usepackage[normalem]{ulem}
\usepackage{hyperref}
\hypersetup{colorlinks=true,linkcolor=black,urlcolor=black,pdfborder={0 0 0}}

% --- Métadonnées du lot -----------------------------------------------------
% Reprises telles quelles par le script de rendu, qui les affiche en tête de
% la vue de lecture. Elles fixent ce que le fichier prétend couvrir : sans
% elles, une transcription flotte sans point d'attache dans un fonds de seize
% mille feuillets.

\newcommand{\folder}[1]{\gdef\grothFolder{#1}}
\newcommand{\batch}[1]{\gdef\grothBatch{#1}}
\newcommand{\leaves}[2]{\gdef\grothFirst{#1}\gdef\grothLast{#2}}
\newcommand{\foldertitle}[1]{\gdef\grothFoldertitle{#1}}
\gdef\grothFolder{?}\gdef\grothBatch{?}\gdef\grothFirst{?}\gdef\grothLast{?}\gdef\grothFoldertitle{}

% --- Le résumé --------------------------------------------------------------
%
% Il ouvre la lecture modernisée, sous le titre « Résumé », et il est composé
% en petit. Ce n'est pas un ornement : c'est le seul passage du document qui
% ne soit pas des mathématiques, et un lecteur doit pouvoir le reconnaître
% avant de l'avoir lu — puis le sauter s'il connaît déjà le sujet, ou s'y
% arrêter s'il ne le connaît pas. Le filet qui le suit marque la reprise du
% corps.

\newenvironment{resume}
  {\small\setlength{\parskip}{0.45em}}
  {\par\medskip\hrule\medskip}

% --- Mots-clés ---------------------------------------------------------------
%
% En anglais, en clôture du résumé de la lecture modernisée : le vocabulaire
% moderne sous lequel chercher ce que les feuillets construisent. Le manifeste
% du site (`npm run manifest`) les extrait pour étiqueter la cote entière —
% cette ligne est la source unique des tags, il n'y a pas de fichier à côté.
\newcommand{\keywords}[1]{\par\smallskip\noindent
  {\footnotesize\textsc{Keywords} --- \textit{#1}}\par}

% --- L'appareil critique ----------------------------------------------------

% Le numéro de feuillet, en marge. C'est l'ancre qui permet de régler un
% désaccord sur une formule en regardant *un* feuillet plutôt qu'une chemise
% de six cents. Le site s'en sert aussi pour tourner les pages du fac-similé
% quand on fait défiler la transcription.
\newcommand{\leaf}[1]{%
  \marginnote{\footnotesize\color{gray!70}\textsf{#1}}%
  \label{leaf:#1}\ignorespaces}

% Illisible : on ne devine pas. Un mot inventé qui se lit comme les autres est
% le pire résultat possible.
\newcommand{\ill}{\textcolor{red!60!black}{[\,\ldots\,]}}

% Lecture douteuse : le mot est donné, et signalé comme incertain.
\newcommand{\uncertain}[1]{\uline{#1}}

% Ajout de l'éditeur — une lettre manquante, un mot sous-entendu.
\newcommand{\add}[1]{\textcolor{blue!50!black}{[#1]}}

% Biffé par Grothendieck lui-même. Ce qu'il a rayé fait partie du document :
% une rature dit souvent où la pensée a changé de direction.
\newcommand{\struck}[1]{\sout{#1}}

% Note du transcripteur, sur l'état matériel du feuillet.
\newcommand{\note}[1]{\par\smallskip{\small\color{gray!60!black}\textit{[#1]}}\par\smallskip}

% Note marginale de Grothendieck, distincte du corps du texte.
\newcommand{\marginal}[1]{\par\smallskip\noindent\hspace{1em}%
  {\small\color{gray!60!black}#1}\par\smallskip}

% --- Titre ------------------------------------------------------------------

\AtBeginDocument{%
  \begin{center}
    {\large\bfseries Fonds Alexandre Grothendieck --- cote n\textsuperscript{o}~\grothFolder}\\[2pt]
    {\itshape\grothFoldertitle}\\[4pt]
    {\small feuillets \grothFirst--\grothLast\ (lot \grothBatch)}\\[2pt]
    {\footnotesize\color{gray!60!black}Transcription établie d'après le fac-similé numérisé
     par l'Université de Montpellier.\\
     Les lectures incertaines sont soulignées, les illisibles notées [\,\ldots\,],
     les ajouts entre crochets.}
  \end{center}
  \vspace{1em}}

\endinput


\folder{12}
\batch{7}
\pages{121}{140}
\dating{1964-[vers 1971]}
\watermark{Édition de démonstration}
\foldertitle{Généralités (sous-groupes de Galois motiviques) (1964 ?) : notes manuscrites (s.d.), tapuscrit (s.d.).}

\begin{document}

\page{121}
2) $F$ \struck{\ill{}} \uncertain{strictement spécial}\add{,} \marginal{\ill{} \uncertain{projectif}} $X$ régulier.
\note{l'insertion interlinéaire au-dessus de « spécial, $X$ » est de sa main ; on n'en lit que le dernier mot, douteux}

\begin{tikzcd}
X' \arrow[r, "f'"] \arrow[d, "\varphi"'] & Y' \arrow[d, "g"] \\
X \arrow[r, "f"'] & \uncertain{Y}
\end{tikzcd}

\note{le carré $X' \to Y' \to$ est entouré d'un trait ; le sommet en bas à droite est écrit en gras, par-dessus une autre lettre, et ne se lit pas sûrement}

\ill{} \ill{} \ill{} $F = R^{d}\varphi_{*}(\mathbb{Q}_{\ell X'})$, \marginal{$\varphi$ de \uncertain{dimension relative} $d$}
\note{la précision « $\varphi$ de dimension relative $d$ » est encadrée}
On a alors
\[
R^{\bullet}(f\varphi)_{*}(\mathbb{Q}_{\ell X'}) \Longleftarrow E_{2}^{pq} = R^{p}f_{*}(R^{q}\varphi_{*}(\mathbb{Q}_{\ell}))
\]
suite spectrale qui \struck{\ill{}} \uncertain{dégénère} par Deligne,
\[
R^{i}f_{*}(F) = R^{i}f_{*}(R^{j}\varphi_{*}(\mathbb{Q}_{\ell X'}))
\]
est un \emph{\ill{}} de $R^{i+j}(f\varphi)_{*}(\mathbb{Q}_{\ell X'})$.
\note{la page écrit tantôt $R^{d}\varphi_{*}$, tantôt $R^{j}\varphi_{*}$ pour le même faisceau $F$ ; transcrit tel quel. Le mot souligné avant « de » se termine peut-être en « -quotient »}

Or aussi :
\[
R^{n}(f\varphi)_{*}(\mathbb{Q}_{\ell X'}) \Longleftarrow E_{2}^{pq} = R^{p}g_{*}(R^{q}f'_{*}(\mathbb{Q}_{\ell X'}))
\]
où $Y'$ est de dim \struck{\ill{}} $n' \leq n+j$, \note{« $n+j$ » est cerclé}
$R^{q}f'_{*}(\mathbb{Q}_{\ell X'})$ est effectif de poids $\leq 2\,\mathrm{Inf}(\ill{})$
\uncertain{au niveau} $\leq \mathrm{Inf}(q, n+j-q)$ ; \ill{} est de \ill{} \uncertain{relatif} $\leq$
\struck{\ill{} $p+q = i+j$} \struck{\ill{}} $E_{2}^{pq}$ de \uncertain{variété} $\leq$ \ill{}

\page{122}
\note{feuillet dactylographié, corrigé de sa main. Ses ajouts manuscrits sont donnés comme notes marginales de sa main ; ce qu'il a biffé à la main, comme ce que la frappe a barré de x, est donné comme biffé}

4) \marginal{\emph{Trivialité cohomologique d'}} \emph{Un morphisme projectif et lisse} $f\colon X \to Y$ \emph{de schémas} (correspondant intuitivement à une fibration sur $Y$ à fibres de variétés projectives non singulières). \struck{est toujours cohomologiquement trivial (i.e.} \marginal{En d'autres termes,} la suite spectrale de Leray pour la cohomologie $\ell$-adique -- ou à coéfficients réels dans le cas du corps de base $\mathbb{C}$ -- dégénère. \struck{\ill{}} \marginal{Ce théorème est prouvé par Deligne modulo le « théorème de Lefschetz » pour la cohomologie des fibres.} Sa démonstration s'applique également \struck{\ill{}} dans le contexte des variétés analytiques, pour une fibration analytique complexe à fibres compactes kählériennes. Un cas très particulier de ce théorème, relatif au cas où le système local de la cohomologie des fibres était trivial, constitu\struck{e}\marginal{ait} le résultat principal de la thèse de Blanchard ; mais l'hypothèse de Blanchard n'est réalisée dans presque aucun cas intéressant en géométrie. La démonstration de Deligne est différente bien entendu de celle de Blanchard, que Deligne ne connaissait pas.
\note{le titre « Trivialité cohomologique d'\,» est ajouté à la main au-dessus de la ligne, et souligné avec « Un morphisme projectif et lisse » et « de schémas » ; le « U » de « Un » est corrigé en minuscule. La ligne dactylographiée qui suit « Sa démonstration s'applique également » est barrée de x à la frappe, et l'ajout « Ce théorème est prouvé par Deligne … » est écrit entre les lignes, au-dessus de cette phrase}

Déjà du point de vue transcendant, l'intérêt du résultat de Deligne pour la topologie des variétés analytiques est évident. Dans l'optique du rédacteur cependant, ce sont ses implications en géométrie algébrique abstraite, et à longue échéance, à la théorie des nombres, qui sont les plus intéressantes, \struck{par exemple, la} \struck{et} Ce sont d'ailleurs des considérations arithmétiques qui avaient amené le \struck{rédacteur} \marginal{rapporteur} à communiquer à Deligne \struck{la conjecture} l'énoncé précédent comme une conjecture. Grâce au théorème de Deligne, on peut voir par exemple que les hypothèses de Riemann-Weil en géométrie diophantienne, interprétées en termes des propriétés d'un endomorphisme de frobénius sur la cohomologie $\ell$-adique, \struck{restent valables sous} impliquent des propriétés analogues dans le cas de coéfficients tordus (pouvant avoir des singularités quelconques, de plus). Le théorème de Deligne sera un des \marginal{nombreux} ingrédients techniques \struck{nécessaires} \marginal{indispensables} pour développer la « théorie des motifs », qui se\ill{}
\note{les mots barrés de x à la frappe sont lus sous la surcharge : « par exemple, la », « la conjecture », « restent valables sous » ; « et », barré à la main, est remplacé par une majuscule à « Ce » ; la fin de la page, « qui se », est surchargée à la main et la phrase se poursuit sur un feuillet qui n'est pas dans ce lot}

\section{\uncertain{Divers types} de $\mathbb{Q}_{\ell}$-faisceaux}
\note{titre souligné, en tête du feuillet, de sa main ; le premier mot est douteux}

\page{124}
\struck{Un faisceau}

Un $\mathbb{Q}_{\ell}$-faisceau $F$ \struck{\ill{}} sur $X$ est dit \marginal{de type $\leq d$} \emph{admissible} s'il satisfait la condition suivante :

$\forall x \in X$, existe voisinage ouvert $U$ de $x$ dans $X$ et filtration \marginal{finie} de $F|U$ par des sous-faisceaux \struck{\ill{}} tels que les \marginal{en.\ constituants} $\mathrm{gr}^{i}(F|U)$ soient isomorphes à des \marginal{\uncertain{sommes directes} d'un} faisceau\supplied{x} de la forme
\[
R^{n}g_{*}(\mathbb{Q}_{\ell}),
\]
où $g_{i}\colon V \to U$ est un morphisme projectif, \uncertain{et} lisse, à fibres connexes \marginal{\uncertain{de dim} $\leq d$}.
\note{dans la parenthèse de $R^{n}g_{*}(-)$ un symbole est noirci avant $\mathbb{Q}_{\ell}$. Une note oblique de sa main dans la marge de gauche, près de cette formule, se lit en partie : « \ill{} en.\ constituants \ill{} faisceaux \ill{} »}

[Se vérifie \uncertain{fibre par fibre}, $+$ $F$ est « constructible »]
\struck{Ceci est \ill{} que $F$ \ill{} est \ill{} admissible} \marginal{\ill{} \uncertain{fibre par fibre}}

\emph{Lemme} 1. Les $\mathbb{Q}_{\ell}$-faisceaux admissibles de type $\leq d$ forment une sous-catégorie \marginal{épaisse} \struck{\ill{} et épaisse (stable par} abélienne \struck{\ill{}} de la catégorie des \struck{$\mathbb{Q}_{\ell}$-faisceaux constructibles (\ill{} noyaux, conoyaux, extensions) de la catégorie des}
\note{la fin de l'énoncé est un enchevêtrement de mots biffés et de mots ajoutés entre les lignes ; on n'en donne que ce qui se lit}

\struck{\ill{}} \uncertain{Elle} est stable par ss-\struck{faisc.} objet et objet quotient. En effet, si $F|U$ est muni d'une filtration \marginal{\uncertain{et si} $G|U$ est en.\ \uncertain{constituants},} de type spécial, \uncertain{alors} la filtration induite sur \ill{} \ill{} de spécial, \struck{\ill{}} \ill{}

\ill{} \ill{} \uncertain{stable par} extensions : trivial.

\emph{Lemme} 2. Image inverse : trivial.
\note{les trois dernières lignes avant le lemme 2 sont d'une écriture très rapide ; seuls les mots ci-dessus se lisent}

\page{126}
\emph{Lemme} \struck{2} 3. \uncertain{Soit} $f\colon X \to Y$ un \marginal{compactifiable} morphisme de type fini, de dim.\ relative $\leq d'$, $F$ un $\mathbb{Q}_{\ell}$-faisceau admissible de type $\leq d$ sur $X$, alors pour tout $i$, $R^{i}_{!}f(F)$ sur $Y$ est admissible de type $\leq d+d'$.
\note{le lemme est marqué dans la marge de gauche par un long trait vertical terminé en flèche. Le « $\mathbb{Q}_{\ell}$ » de « $\mathbb{Q}_{\ell}$-faisceau » ressemble ici à un $\mathbb{Z}_{\ell}$ ; les autres occurrences de la page et de la précédente sont des $\mathbb{Q}_{\ell}$. Il écrit $R^{i}_{!}f$ pour $R^{i}f_{!}$}

On peut supposer $Y = \operatorname{Spec} k$, $k$ un corps.

Récurrence sur $d+d'$, \uncertain{trivial} si $d' \leq 0$. Supposons $d+d' > 0$ et le th.\ démontré pour les \struck{\ill{}} $\delta' < d'$.

Soit $U \subset X$ \marginal{\uncertain{ouvert}} lisse sur $K$ et dense dans $X$, et tel que $F|U$ ait une filtration d'admissibilité.
[N.B. \uncertain{On} \uncertain{peut} \uncertain{exiger} dans la déf.\ plus haut \uncertain{d'admettre} une extension \uncertain{radicielle} sur la base…]
\note{« $\delta' < d'$ » est écrit au-dessus d'un mot biffé. Dans le N.B., le verbe « exiger » et le mot qui précède « sur la base » sont douteux ; les points de suspension sont de sa main}

Soit $Z = X - U$. On a
\[
R^{i}_{!}f_{U}(F|U) \to R^{i}_{!}f_{*}(F) \to R^{i}_{!}f_{Z*}(F|Z)
\]
ce qui nous ramène à prouver le th.\ pour $f_{Z}$ et $f_{U}$ grâce aux lemmes 1 et 2 ; or le cas de $f_{Z}$ est couvert par l'hyp.\ de réc., et \struck{\ill{}} finalement on est ramené au cas où $X$ est lisse sur $k$ et a une filtration d'admissible. Par \uncertain{dévissage} et le lemme 1 on peut supposer que $F$ est \marginal{en.\ const.} un \uncertain{quotient} d'un
\note{la phrase s'arrête ici, au bas du feuillet ; elle reprend au feuillet 128}

\page{128}
\note{suite de la phrase interrompue au bas du feuillet 126}
somme directe \marginal{en.\ const.} d'un \ill{} $R^{m}g_{*}(\mathbb{Q}_{\ell})$, \ill{} $g\colon V \to U$ un morphisme \struck{projectif} de \marginal{relative} dim.\ $\leq d$. On a \ill{} \ill{} \ill{}, \uncertain{des} \uncertain{suites} \uncertain{exactes}
\[
0 \to G \to R^{m}g_{*}(\mathbb{Q}_{\ell}) \to H \to 0
\]
\[
0 \to I \to G \to F \to 0
\]
\note{un astérisque de sa main est posé sous $I$, sous $G$ et sous $F$ dans la seconde suite ; un trait oblique sous le $G$ de la première. Le « \uncertain{ex. dans} » écrit au-dessus, à gauche des suites, ne se lit pas sûrement}

et de nouvelles applications du lemme 2 nous ramènent au cas de $I$ et $G$, $R^{n}f_{*}(F)$ lorsque \struck{$I$, $G$} $F$ en.\ const.\ $= R^{m}g_{*}(\mathbb{Q}_{\ell})$. Pour conclure, il faut utiliser le « \uncertain{théorème} de \ill{} » \ill{} complète admissibilité.
\note{la fin de la phrase est soulignée d'un trait qui court sous toute la ligne}

\uncertain{Pour} \ill{}, \uncertain{soumis} \ill{} \ill{} \ill{} \ill{}. $F$ \uncertain{est} de la forme $R^{m}g_{*}(\mathbb{Q}_{\ell})$. \uncertain{Reste} deux cas à examiner. Ce cas :
\[
R^{p}_{!}f_{*}(R^{q}g_{*}(\mathbb{Q}_{\ell})) \qquad p, q \geq 0
\]

\struck{b) \emph{Cas $m = 0$}. Nous procédons par réc.\ \emph{descendante} sur $p+q$ \ill{}, \uncertain{on} \uncertain{suppose} \ill{} \uncertain{vrai} pour les $(p', q')$ avec $p' > p$ \ill{}}
\note{ces trois dernières lignes sont encadrées à gauche d'un crochet et barrées de trois longues diagonales}

\page{130}
De plus, procédant par récurrence sur $d+d'$, et en utilisant une \emph{factorisation d'Artin} de $U$ [ce qui est loisible, car on \uncertain{peut} \uncertain{prendre} $U$ \emph{petit}], on \uncertain{peut} \uncertain{supposer} \uncertain{que} $U$ est de \emph{dim} 1.

Alors la suite \uncertain{exacte} \ill{} \ill{} $g$, \uncertain{compte} \uncertain{tenu} de
\[
R^{0}_{!}f(R^{q}g_{*}(\mathbb{Q}_{\ell})) = 0,
\]
se \uncertain{réduit} à \struck{\ill{}} \uncertain{des} suites exactes courtes
\[
0 \to R^{2}_{!}f_{*}(R^{n-2}g_{*}(\mathbb{Q}_{\ell})) \to R^{n}_{!}(gf)_{*}(\mathbb{Q}_{\ell}) \to R^{1}_{!}f_{*}(R^{n-1}g_{*}(\mathbb{Q}_{\ell})) \to 0
\]
\note{le $0$ final est écrit sous le dernier terme, au bout d'une petite flèche verticale}

donc par le lemme 1 il suffit de vérifier que
\[
R^{n}_{!}(gf)_{!}(\mathbb{Q}_{\ell})
\]
est spécial. \note{« $h$ » est écrit au-dessus de $gf$ : c'est la lettre qui sert dans la suite}
Or utilisons la résolution des singularités pour $V$,
\[
V = \hat{V} - Z, \qquad \hat{V} \text{ projectif et lisse sur } k,
\]
d'où \uncertain{la} \uncertain{suite} \ill{}
\[
R^{n-1}_{!}h_{Z}(\mathbb{Q}_{\ell}) \to R^{n}_{!}h(\mathbb{Q}_{\ell}) \to R^{n}_{!}\hat{h}(\mathbb{Q}_{\ell})
\]
\note{le terme du milieu porte au-dessus du $h$ un signe (chapeau ou astérisque) qui ne se distingue pas sûrement de celui du troisième terme ; la page s'arrête sur cette suite}

\page{132}
et \struck{\ill{}} par le lemme 1 il suffit \ill{} \ill{} \ill{} les \ill{} \emph{catégories}. Celui-ci a \ill{} \ill{} dont \ill{} est O.K. \struck{\ill{}} par \emph{définition}, \ill{} \ill{} par hyp.\ de récurrence. O.K.
\note{le haut du feuillet, qui achève sans doute un raisonnement commencé ailleurs, est d'une écriture très rapide ; seuls les mots ci-dessus se lisent}

\emph{Proposition} 1. \uncertain{Soit} $i\colon U \to X$ une immersion ouverte, avec $X$ excellent, de dim $\leq d$, $F$ un $\mathbb{Q}_{\ell}$-faisceau \struck{\ill{}} $d'$-admissible. Alors $R^{n}i_{*}(F)$ est $(d+d'\uncertain{-1})$-admissible.
\note{l'énoncé est marqué d'un trait vertical dans la marge. « avec $X$ excellent, de dim $\leq d$ » est ajouté entre les lignes. Dans l'indice final, le dernier chiffre est biffé et surchargé d'un « 1 » écrit au-dessus ; la lecture « $d+d'-1$ » est douteuse}

\struck{admissible.}

\struck{En effet, \ill{} $V \subset U$, $W$ \ill{}, $W$ \ill{} \ill{} $U$, tel \ill{}}
\note{deux lignes barrées de longs traits ondulés}

Récurrence sur $d+d'$, trivial si $d+d' = 0$ i.e.\ $d = 0$, $d' = 0$. Supposons $d+d' > 0$, et le th.\ \uncertain{prouvé} si $\delta + \delta' < d+d'$. \ill{} \uncertain{peut} \uncertain{supposer} $X$ \ill{} \ill{} \ill{}. \struck{\ill{}} Le th.\ est \ill{} \uncertain{par} \uncertain{récurrence} \uncertain{noethérienne} \uncertain{sur} les $Y \subset X$ fermés, $Y \neq X$, et $X \neq \emptyset$.

Soit donc $V \subset U$ \ill{} dense de $U$, régulier, tel que
\[
F \to \mathbf{R}j_{*}F \to Q^{*},
\]
où $Q^{*}$ est concentré sur $Z \cap U$, $Z = \overline{U - V}$ \uncertain{vu} dans $X$, \ill{} $\neq X$. Donc \ill{} \ill{} \ill{} \ill{} prouver que les $\mathbf{R}^{n}i_{*}(\mathbf{R}j_{*}(F))$ et $\mathbf{R}^{n}i_{*}(Q^{*})$ sont admissibles. Par la
\note{$j\colon V \to U$ ; $Q^{*}$ est écrit au-dessus d'une flèche partant de $\mathbf{R}j_{*}F$, comme troisième sommet du triangle. La phrase se poursuit au feuillet 134}

\page{134}
\note{suite de la phrase interrompue au bas du feuillet 132}
lemme 1 et l'hyp.\ de récurrence, \uncertain{c'est} OK pour $\mathbf{R}^{n}i_{*}(Q^{*})$. Reste : $\mathbf{R}^{n}i_{*}(\mathbf{R}j_{*}(F))$ i.e.\ aussi $\mathbf{R}^{n}(ij)_{*}(F)$, ce qui nous ramène au cas où $U$ est \emph{régulier}. Alors par \emph{résolution} \ill{}, $i = fi'$, avec

\begin{tikzcd}
& X' \arrow[d, "f"] \\
U \arrow[ur, hook, "i'"] \arrow[r, hook, "i"'] & X
\end{tikzcd}

$f$ \uncertain{projectif}, $X'$ régulier, $i'$ une immersion \ill{} à complémentaire un diviseur à croisements normaux.

Mais
\[
\mathbf{R}^{*}\uncertain{f_{*}}(F) \Longleftarrow \mathbf{R}^{p}f_{*}(R^{q}i'_{*}(F))
\]
et cela nous ramène par le lemme 3 au cas où \struck{$X$} $X - U$ est un diviseur à croisements normaux. \emph{Lorsque} $F = \mathbb{Q}_{\ell}$, on \uncertain{gagne} alors par (th.\ \emph{de pureté}). [On va ramener le cas général à celui-ci, \ill{} \ill{} : celui-ci, on compactifie $V$ \struck{\ill{}} en \struck{\ill{}} $g\colon Y \to X$ morphisme projectif (par \ill{} \ill{}), où $F$ \ill{} \marginal{\ill{} de dim.\ rel.\ $g \leq d'$} facteur direct de $R^{n}t_{*}(\mathbb{Q}_{\ell})$ ; (on a

\begin{tikzcd}
V \arrow[r, hook, "j"] \arrow[d, "t"'] & \hat{Y} \arrow[d, "g"] \\
U \arrow[r, hook, "i"'] & X
\end{tikzcd}

\note{le premier membre de la suite spectrale est d'une lecture douteuse : on attendrait $R^{*}(fi')_{*} = R^{*}i_{*}$ ; transcrit tel qu'il paraît. Au-dessus de $Y$, dans le carré, un accent circonflexe, que le texte n'a pas}

Posons $F = R^{n}t_{*}(\mathbb{Q}_{\ell}) = R^{n}j_{*}(\mathbb{Q}_{\ell})|U$. Par \ill{} des diviseurs \ill{} ailleurs, utilisant \marginal{lemme 3} \marginal{\ill{} \ill{} \ill{} \ill{}}, on se ramène au cas où $d' = 1$. Alors le \uncertain{cas} spécial \ill{} \ill{} \ill{}
\note{dans la marge de gauche, en bas, une note oblique soulignée, de sa main, dont on ne lit rien sûrement. La page s'arrête ici}

\page{136}
\[
\mathbf{R}^{*}(if)_{*} \Longleftarrow \mathbf{R}^{p}i_{*}(\mathbf{R}^{q}f_{*}(\mathbb{Q}_{\ell}))
\]
\ill{} \ill{} $E_{2}^{pq} = 0$ si $q \notin [0, 2]$.
\note{sous la formule, un croquis du terme $E_{2}$ : trois lignes horizontales, la plus haute marquée « 2 », fermées à droite par un trait vertical marqué « $N$ » ; des points sur les lignes, et deux flèches obliques $d_{2}$ descendant de la ligne 2 vers la ligne 1, puis de la ligne 1 vers la ligne 0, vers la droite}

\uncertain{(Car} \ill{} \ill{} si $f$ \ill{} \ill{} \ill{} \uncertain{revient} \ill{} \ill{} \ill{}.) D'ailleurs $\mathbf{R}^{0}f_{*}(\mathbb{Q}_{\ell})$ est essentiellement $\mathbb{Q}_{\ell}$ [\ill{} \uncertain{à} fibres connexes, \ill{} \ill{} \uncertain{facteurs} …], donc les \struck{$R^{p}i_{*}$ \ill{}} $E_{2}^{p0}$ sont $(d-1)$-admissibles \ill{} \ill{} \uncertain{hyp.} \uncertain{récurr.}
\note{les points de suspension entre crochets sont de sa main}

\struck{Donc} D'autre part $R^{2}f_{*}(\mathbb{Q}_{\ell}) \simeq \mathbb{Q}_{\ell}(-1)$, et
\[
E_{2}^{p2} \simeq \mathbf{R}^{p}i_{*}(\mathbf{R}^{2}f_{*}(\mathbb{Q}_{\ell})) \simeq \mathbb{Q}_{\ell}(-1) \otimes E_{2}^{p0}
\]
est donc \uncertain{également} $(d-1)$ \uncertain{admissible}, \uncertain{i.e.} \uncertain{fortiori} \marginal{\ill{} $d$-admissible,} $(d+d'-1)$-admissible. \struck{Pour} \uncertain{Enfin} \uncertain{on} \uncertain{voit} \marginal{et la \uncertain{constructibilité} des $E_{2}^{pq}$} \ill{} utilisant lemme 2, que $E_{2}^{p1}$ est une extension de \struck{\ill{}} \marginal{trois} faisceaux, les \uncertain{autres} \struck{$(d+d'-1)$ \ill{}} $d$-admissible, et le \uncertain{troisième} \ill{} $E_{3}^{p1} = \cdots = E_{\infty}^{p1}$, \ill{} \uncertain{effet} \ill{} \ill{} prouver que $E_{\infty}^{p1}$ est $d$-admissible, \ill{} \ill{} \ill{} que $\mathbf{R}^{p+1}(gf)_{*}(\mathbb{Q}_{\ell}|V)$, \ill{} \ill{}
\ill{} pour $R^{n}(gf)_{*}(\ill{})$ \ill{} \ill{} \ill{} \ill{} \ill{} \uncertain{lemme} et lemme 3 \ill{}, \uncertain{ramène} au \uncertain{cas} \ill{} \ill{} $R^{n}j_{*}(\mathbb{Q}_{\ell})$, qui \uncertain{reste} \uncertain{à} \uncertain{traiter}…!!
\note{la seconde moitié de la page est d'une écriture très rapide, chargée d'ajouts interlinéaires ; les formules se lisent, la prose presque pas. Les points de suspension et le double point d'exclamation final sont de sa main}

\section{Motifs lisses versus Isomotifs lisses}
\note{titre encadré, seul sur le feuillet 138, de sa main}

\page{139}
$S$ schéma normal intègre, pt générique $\eta$, $k(\eta) = K$.

\emph{Isomotif} ($\tfrac{1}{2}$ simple) \emph{lisse} sur $S$ : isomotif semi-simple $M$ \marginal{$M$} sur $S$, tel que pour tout $\ell \in \mathbb{P}$, $\ell \neq \mathrm{car}\, K$, $T_{\ell}(M)$ soit non ramifié rel.\ à $S_{\ell} = S[1/\ell]$.
Conjecturalement, il suffit que ce soit vrai pour \emph{un} $\ell$ premier à toutes les car.\ résiduelles ($\Longrightarrow$ il suffit que ce soit vrai pour \emph{deux} $\ell \neq \ell'$).
\note{« $\tfrac{1}{2}$ simple » : sa notation de « semi-simple », qu'il écrit en toutes lettres juste après. Le « $M$ » est ajouté au-dessus de la ligne, avec un trait d'insertion avant « sur $S$ ». L'isomotif étant pris sur le corps $K$, on attendrait « sur $K$ » là où la page écrit « sur $S$ »}

Motif (\struck{\ill{}} isosemisimple) lisse sur $S$ (si $\mathrm{car}\, K = 0$) :
\[
= \bigl\{ M, \{E_{\ell}\}, \{\alpha_{\ell}\} \bigr\}
\]
où
\begin{enumerate}
  \item[a)] $M$ un isomotif lisse sur $S$ ;
  \item[b)] $\forall \ell$, $E_{\ell}$ un faisceau $\ell$-adique et tordu sur $S_{\ell} = S[1/\ell]$ [ou encore sur $K$, mais non ramifié rel.\ à $S[1/\ell]$] ;
  \item[c)] $\forall \ell$, $\alpha_{\ell}\colon E_{\ell} \otimes_{\mathbb{Z}_{\ell}} \mathbb{Q}_{\ell} \xrightarrow{\ \sim\ } T_{\ell}(M)$,
\end{enumerate}
avec un axiome : « pour presque tout $\ell$, $u_{\ell}$ provient d'un isomorphisme sur $\mathbb{Z}_{\ell}$ -- pas seulement : torsion près ».
\note{le produit tensoriel de c) est surchargé ; « et tordu » en b) est une lecture douteuse. Dans l'axiome, il écrit $u_{\ell}$ là où l'on attend $\alpha_{\ell}$}

Ceci indique que pour \struck{presque} tout $\ell$, $E_{\ell}$ est sans torsion.

Homomorphismes de motifs lisses sur $S$
\[
\bigl(M, \{E_{\ell}\}, \{\alpha_{\ell}\}\bigr) \to \bigl(M', \{E'_{\ell}\}, \{\alpha'_{\ell}\}\bigr)
\]
données par $(u, \{u_{\ell}\})$, avec $u\colon M \to M'$ un hom.\ de \uncertain{isomotifs}, et $u_{\ell}\colon E_{\ell} \to E'_{\ell}$ un hom.\ de faisceaux $\ell$-adiques tels que pour tout $\ell$, \uncertain{on} ait commutativité \uncertain{dans}

\begin{tikzcd}
E_{\ell} \otimes_{\mathbb{Z}_{\ell}} \mathbb{Q}_{\ell} \arrow[r, "\alpha_{\ell}"] \arrow[d, "u_{\ell} \otimes_{\mathbb{Z}_{\ell}} \mathbb{Q}_{\ell}"'] & T_{\ell}(M) \arrow[d, "T_{\ell}(u)"] \\
E'_{\ell} \otimes_{\mathbb{Z}_{\ell}} \mathbb{Q}_{\ell} \arrow[r, "\alpha'_{\ell}"'] & T_{\ell}(M')
\end{tikzcd}

\emph{N.B.} Le foncteur \struck{\ill{}} (Motifs sur $K$) $\longleftarrow$ (Motifs \uncertain{\struck{lisses}} sur $S$) est pleinement fidèle.
\note{« Motifs » est écrit sous un mot biffé, à gauche ; à droite, le mot qui suit « Motifs » paraît biffé, sans certitude} Les Hom sont des \uncertain{modules} de type fini sur $\mathbb{Z}$. La catégorie des motifs lisses sur $S$ est abélienne (on prend \struck{\ill{}} noyaux \uncertain{et} \uncertain{conoyaux} \ill{} \ill{} $M$, $E_{\ell}$). \struck{\ill{}} Le composé \marginal{lisses} (Motifs lisses sur $S$) $\to$ (Motifs sur $K$)
\struck{\ill{} \ill{} Mot} (Motifs \uncertain{lisses} \ill{}) est \uncertain{exact} \uncertain{et} \uncertain{fidèle}, identifiant \uncertain{à} une sous-catégorie épaisse. \uncertain{Le} \uncertain{foncteur} \marginal{lisses} (Motifs) $\to$ (Isomotifs) \marginal{lisses} est exact \uncertain{essentiellement} surjectif
\note{le bas du feuillet est une suite d'ajouts entre les lignes et de mots biffés ; l'ordre de lecture proposé ici est incertain}

\page{140}
La catégorie des groupes finis \marginal{commut.} \uncertain{$G$} finis sur $K$ tels que pour tout $\ell$, $G(\ell)$ soit non ramifié rel.\ à $S[1/\ell]$, est une sous-catégorie épaisse de celle des motifs sur $S$, la catégorie quotient est équivalente : celle des isomotifs lisses sur $S$.
\note{« commut.\ $G$ » est ajouté au-dessus de la ligne, entre « finis » et « sur $K$ », et le premier « finis » semble repris par le second ; lecture de l'ensemble douteuse. $G(\ell)$ : sa composante $\ell$-primaire, semble-t-il}

La catégorie des motifs est une $\otimes$-catégorie avec \uncertain{$\underline{\mathrm{Hom}}$}. Les foncteurs Motifs $\to$ Isomotifs, Motifs $\to$ faisceaux $\ell$-adiques, sont des $\otimes$-foncteurs, compatibles aux \uncertain{$\underline{\mathrm{Hom}}$}.
\note{les deux « Hom » soulignés sont écrits seuls sur leur ligne, au bout de chacune des deux phrases}

\struck{Une définition} \ill{} \ill{} \ill{} \ill{} un isom.\ sur les isomotifs \uncertain{sous-jacents}. \uncertain{En} \uncertain{effet} \ill{} la dimension (\uncertain{hors} le \ill{} \uncertain{fibres} \uncertain{vides}), \struck{la} \uncertain{niveau} etc.\ en termes des isomotifs.

Tout schéma projectif et lisse $X$ sur $S$ définit des motifs \marginal{sur $S$}, notés $R^{i}f_{*}(\mathbb{Z}_{X})$ s'il \uncertain{n'y} a pas de confusion à craindre. \ill{} \ill{} \ill{} \ill{} \uncertain{motifs} \struck{lisses} et \uncertain{propres}.

\emph{Relations avec les schémas en groupes commutatifs propres et lisses.}

On définit un foncteur \emph{contravariant} de la catégorie des \marginal{schémas en groupes} \uncertain{A} \uncertain{lisses} \marginal{sur $S$}, dans la catégorie des motifs sur $S$, en posant
\struck{$\mathfrak{P}(A) = R^{1}f_{*}(\mathbb{Z}_{A}) = \{R^{1}f_{*}($}

\[
\mathfrak{P}(A) = (M, E_{\ell}, u)
\]
où
\[
\begin{cases}
M = R^{1}(f_{A^{0}})_{*}(\mathbb{Q}_{A^{0}}) \\
E_{\ell} = \underline{\mathrm{Hom}}({}_{\ell^{\infty}}A, \mathbb{Q}_{\ell}/\mathbb{Z}_{\ell}) \quad \text{extension de } \underline{\mathrm{Hom}}(A/A^{0}, \mathbb{Q}_{\ell}/\mathbb{Z}_{\ell}) \text{ par } R^{1}f_{A^{0}*}(\mathbb{Z}_{\ell A^{0}}) \\
u_{\ell} \text{ provenant de l'hom.\ induit } E_{\ell} \to R^{1}f_{A^{0}*}(\mathbb{Z}_{\ell A^{0}})
\end{cases}
\]
\note{la première ligne de définition, $\mathfrak{P}(A) = R^{1}f_{*}(\mathbb{Z}_{A}) = \{R^{1}f_{*}(\,$, est barrée de deux traits et reprise sur la ligne suivante. Le $\mathfrak{P}$ est une lecture : une capitale gothique ou ronde. ${}_{\ell^{\infty}}A$ : la torsion $\ell$-primaire de $A$, l'indice écrit à gauche}

\emph{NB} Il suffirait que $A$ \struck{\ill{} extension \ill{} \ill{}} \ill{} \ill{} $A/A^{0}$ \ill{} un schéma abélien $A^{0}$, avec $A/A^{0}(\ell)$ extension \ill{} d'un groupe fini étale sur $S_{\ell}$.
\note{ce N.B. est corrigé en cours d'écriture, et sa construction ne se lit pas sûrement. La page s'arrête ici}

\end{document}
